\begin{figure}[htbp]
\centering
\resizebox{0.92\textwidth}{!}{%
\begin{tikzpicture}[
  font=\sffamily\small,
  block/.style={draw=VIUOrange,line width=0.9pt,rounded corners=4pt,
    align=center,text width=3.05cm,minimum height=1.05cm,inner sep=6pt,fill=VIUOrange!7},
  gate/.style={draw=VIUDark!55,line width=0.8pt,rounded corners=4pt,
    align=center,text width=9.2cm,minimum height=0.88cm,inner sep=6pt,fill=VIUWhite},
  flow/.style={-{Latex[length=2.3mm]},line width=0.9pt,draw=VIUDark!75}
]
\node[block] (g) at (0,1.4) {1. Geometr\'ia\\normal y contacto};
\node[block] (w) at (3.8,1.4) {2. \emph{Wrench}\\error f\'isico};
\node[block] (p) at (7.6,1.4) {3. Fase\\Smith--Kuramoto};
\node[block] (r) at (7.6,-0.35) {4. Rigidez\\marco local};
\node[block] (c) at (3.8,-0.35) {5. Filtro\\CBF--QP};
\node[block] (u) at (0,-0.35) {6. Uniciclo\\ruedas};

\draw[flow] (g) -- (w);
\draw[flow] (w) -- (p);
\draw[flow] (p) -- (r);
\draw[flow] (r) -- (c);
\draw[flow] (c) -- (u);

\node[gate] (gate) at (3.8,-2.05) {Criterios medibles: residual \emph{wrench}, QP factible,\\saturaci\'on de uniciclo, normal finita y eventos sin Zeno};
\draw[flow,dashed] (c.south) -- (gate.north);
\end{tikzpicture}}
\caption{Bucle finito de control vectorial para bajar el marco integrado a una
implementación tipo CoppeliaSim. Cada bloque produce una señal medible y un criterio de
verificación local.}
\viuownsource
\label{fig:control-vectorial-seis-bloques}
\end{figure}
